\begin{figure*}[t!]
    \centering
    \begin{minipage}{0.63\linewidth} % Adjust width as needed
        \centering
        \includegraphics[width=0.9\linewidth]{Figures/temporal_axis_scaling/First_occurence_of_higher_reward.PNG}
    \end{minipage}%
    \hfill
    \begin{minipage}{0.35\linewidth} % Adjust caption width
        \scriptsize
        \captionof{figure}{\textbf{Analysis of number of function evaluations (NFEs) across timesteps.}  
        The NFEs required to achieve a higher reward for each timestep. The plot illustrates the $\pm 1$ sigma variation band. The blue-dotted line represents the uniform allocation of compute (NFEs) across timesteps. We observe that the NFEs required to identify a higher-reward sample may exceed the uniformly allocated budget (blue dotted line).}
        \label{fig:temporal_axis_scaling}
    \end{minipage}
\end{figure*}


% \begin{figure}[t!]
% {\scriptsize
% \setlength{\tabcolsep}{0.0em}
% \def\arraystretch{0.0}
% \newcolumntype{Z}{>{\centering\arraybackslash}m{\textwidth}}
%     \begin{tabularx}{\textwidth}{Z}
%       \includegraphics[width=\textwidth]{Figures/temporal_axis_scaling/First_occurence_of_higher_reward.PNG}
%     \end{tabularx}
%   \caption{\textbf{Analysis of number of function evaluations (NFEs)} The NFEs required to achieve a higher reward. The plot illustrates the $±1$ sigma variation band.}
%   \label{fig:temporal_axis_scaling}
%   \vspace{-2.0\baselineskip}
% }
% \end{figure}

